\chapter{Összefoglalás}

A dolgozat egy Java alapú e-learning webalkalmazás — a ME Learning platform — tervezését és megvalósítását mutatta be. Az alkalmazás Spring Boot 3.x keretrendszerre épül, Thymeleaf sablon-motorral és H2 memória-adatbázissal, és a Miskolci Egyetem hallgatói számára nyújt kurzuskezelési, prezentáció-megjelenítési és kvíz-alapú tudásellenőrzési funkciókat.

A fejlesztés során a következő főbb eredmények születtek:

\begin{itemize}
    \item Szerepkör-alapú hozzáférés-vezérlés megvalósítása három szerepkörrel (adminisztrátor, oktató, hallgató), Spring Security segítségével.
    \item Kurzuskezelő rendszer, amely lehetővé teszi kurzusok létrehozását, módosítását, hallgatók beiratkoztatását és eltávolítását.
    \item Prezentáció-feldolgozó modul, amely PowerPoint fájlokat (.ppt, .pptx) automatikusan PDF formátumba konvertál, magas felbontásban renderelve az egyes diákat az Apache POI és PDFBox könyvtárak segítségével.
    \item Kvízrendszer, amely az oktatók által létrehozott teszteket kezeli, eredményeket rögzít, statisztikákat számít, és visszamenőleges osztályzat-újraszámítást tesz lehetővé.
    \item E-mail értesítési funkció Gmail SMTP-n keresztül.
    \item Reszponzív, felhasználóbarát felhasználói felület Thymeleaf sablonokkal és egyedi CSS stíluslapokkal.
\end{itemize}

Az alkalmazás a kitűzött funkcionális követelményeket teljesíti. A manuális tesztelés során az összes azonosított teszteset sikeresen lefutott. A rendszer könnyen bővíthető: a H2 adatbázis éles környezetben PostgreSQL-re cserélhető, és az architektúra lehetővé teszi REST API réteg hozzáadását is.

A fejlesztés során szerzett tapasztalatok alapján a jövőbeli fejlesztési irányok között szerepel a fájlok adatbázison kívüli tárolása, a részletesebb teljesítménystatisztikák bevezetése, valamint a mobilalkalmazás-integráció lehetősége.
